\chapter{Технологическая часть}

В данном разделе представлены выбранные средства реализации, примеры запуска реализованного HTTP-сервера и выполнения функциональных тестов.

\section{Выбор средств реализации}

В качестве целевой операционной системы была выбрана Ubuntu 22.04 LTS~\cite{ubuntu_2204_help}, т.~к. данная операционная система предоставляет поддержку необходимых функций ядра, описанных стандартом POSIX, а также в связи с наличием утилиты wrk~\cite{ubuntu_wrk}, используемой для нагрузочного тестирования HTTP-сервера.

Для реализации HTTP-сервера был выбран язык программирования C~\cite{iso9899_2011}, т.~к. данный язык обеспечивает доступ к необходимым системным вызовам посредством библиотеки glibc~\cite{glibc_manual}.

\section{Реализация алгоритмов}

Ниже приведены листинги основных алгоритмов HTTP-сервера:
\begin{itemize}
	\item листинг~\ref{lst:main}~---~алгоритм работы процесса-предка;
	\item листинг~\ref{lst:signals}~---~обработчики сигналов и их установка;
	\item листинг~\ref{lst:createWorkers}~---~создание процессов-обработчиков;
	\item листинг~\ref{lst:listen}~---~создание слушающего сокета;
	\item листинг~\ref{lst:workerLoop}~---~алгоритм работы worker'а;
	\item листинг~\ref{lst:connection}~---~структура соединения;
	\item листинг~\ref{lst:httpResponse}~---~структура плана HTTP-ответа;
	\item листинг~\ref{lst:httpResponseResult}~---~структура результата отправки HTTP-ответа;
	\item листинг~\ref{lst:responsePlan}~---~функция формирования плана HTTP-ответа;
	\item листинг~\ref{lst:responsePlanFile}~---~функция подготовки файла для HTTP-ответа;
	\item листинг~\ref{lst:make}~---~makefile для сборки и запуска сервера.
\end{itemize}

\begin{lstlisting}[caption={Алгоритм работы процесса-предка}, label={lst:main}]
#define DEFAULT_PORT 8080
#define DEFAULT_WORKERS 4
#define LISTEN_BACKLOG 128
#define DEFAULT_DOCROOT "./www"
#define DEFAULT_LOG_PATH "./server.log"
#define DEFAULT_ERROR_LOG_PATH "./server.error.log"
int main(int argc, char **argv) {
	int port = DEFAULT_PORT;
	int worker_count = DEFAULT_WORKERS;
	const char *docroot = DEFAULT_DOCROOT;
	const char *log_path = DEFAULT_LOG_PATH;
	const char *error_log_path = DEFAULT_ERROR_LOG_PATH;
	for (int i = 1; i < argc; ++i) {
		if (strcmp(argv[i], "-p") == 0 && i + 1 < argc) {
			if (parse_positive_number(argv[++i], &port) < 0 || port > 65535) {
				fprintf(stderr, "Invalid port value\n");
				return 1;
			}
		} else if (strcmp(argv[i], "-w") == 0 && i + 1 < argc) {
			if (parse_positive_number(argv[++i], &worker_count) < 0) {
				fprintf(stderr, "Invalid workers value\n");
				return 1;
			}
		} else if (strcmp(argv[i], "-d") == 0 && i + 1 < argc) {
			docroot = argv[++i];
		} else if (strcmp(argv[i], "-l") == 0 && i + 1 < argc) {
			log_path = argv[++i];
		} else if (strcmp(argv[i], "-e") == 0 && i + 1 < argc) {
			error_log_path = argv[++i];
		} else {
			fprintf(stderr,
			"Usage: %s [-p port] [-w workers] [-d docroot] [-l log_path] "
			"[-e error_log_path]\n",
			argv[0]);
			return 1;
		}
	}
	if (logger_set_error_path(error_log_path) < 0) {
		fprintf(stderr, "Invalid error log path\n");
		return 1;
	}
	if (logger_set_access_path(log_path) < 0) {
		logger_log_error("master", "invalid access log path", EINVAL);
		fprintf(stderr, "Invalid log path\n");
		return 1;
	}
	if (worker_set_docroot(docroot) < 0) {
		logger_log_error("master", "invalid docroot", EINVAL);
		fprintf(stderr, "Invalid docroot\n");
		return 1;
	}
	if (worker_set_log_path(log_path) < 0) {
		logger_log_error("master", "failed to set access log path in worker", EINVAL);
		fprintf(stderr, "Invalid log path\n");
		return 1;
	}
	if (worker_set_error_log_path(error_log_path) < 0) {
		logger_log_error("master", "failed to set error log path in worker", EINVAL);
		fprintf(stderr, "Invalid error log path\n");
		return 1;
	}
	int listen_fd = create_listen_socket(port);
	if (listen_fd < 0) {
		return 1;
	}
	pid_t *workers = calloc((size_t)worker_count, sizeof(pid_t));
	if (workers == NULL) {
		logger_log_error("master", "calloc failed", errno);
		close(listen_fd);
		return 1;
	}
	install_signal_handlers();
	printf("[master %d] listening on port %d, workers=%d, docroot=%s, log=%s, error_log=%s\n",
	getpid(), port, worker_count, docroot, log_path, error_log_path);
	fflush(stdout);
	for (int i = 0; i < worker_count; ++i) {
		pid_t pid = create_worker(listen_fd);
		if (pid < 0) {
			g_stop = 1;
			worker_count = i;
			break;
		}
		workers[i] = pid;
	}
	while (!g_stop) {
		pause();
		if (g_child_changed) {
			g_child_changed = 0;
			while (1) {
				int status = 0;
				pid_t dead = waitpid(-1, &status, WNOHANG);
				if (dead <= 0) {
					break;
				}			
				int slot = find_worker_slot(workers, worker_count, dead);
				if (slot >= 0) {
					workers[slot] = -1;
					if (!g_stop) {
						pid_t replacement = create_worker(listen_fd);
						if (replacement > 0) {
							workers[slot] = replacement;
							printf("[master] respawn worker: %d -> %d\n", dead, replacement);
							fflush(stdout);
						} else {
							logger_log_error("master", "failed to respawn worker", errno);
							g_stop = 1;
						}
					}
				}
			}
		}
	}
	printf("[master] stopping\n");
	fflush(stdout);
	for (int i = 0; i < worker_count; ++i) {
		if (workers[i] > 0) {
			kill(workers[i], SIGTERM);
		}
	}
	while (waitpid(-1, NULL, 0) > 0) {}	
	close(listen_fd);
	free(workers);
	return 0;
}
\end{lstlisting}
\newpage

\begin{lstlisting}[caption={Обработчики сигналов и функция их установки}, label={lst:signals}]
static void on_stop_signal(int signo) {
	(void)signo;
	g_stop = 1;
}
static void on_sigchld(int signo) {
	(void)signo;
	g_child_changed = 1;
}
static void install_signal_handlers(void) {
	struct sigaction sa_stop;
	memset(&sa_stop, 0, sizeof(sa_stop));
	sa_stop.sa_handler = on_stop_signal;
	sigemptyset(&sa_stop.sa_mask);
	sigaction(SIGINT, &sa_stop, NULL);
	sigaction(SIGTERM, &sa_stop, NULL);
	struct sigaction sa_chld;
	memset(&sa_chld, 0, sizeof(sa_chld));
	sa_chld.sa_handler = on_sigchld;
	sigemptyset(&sa_chld.sa_mask);
	sa_chld.sa_flags = SA_RESTART | SA_NOCLDSTOP;
	sigaction(SIGCHLD, &sa_chld, NULL);
}
\end{lstlisting}

\begin{lstlisting}[caption={Функция создания процессов-обработчиков}, label={lst:createWorkers}]
static pid_t create_worker(int listen_fd) {
	pid_t pid = fork();
	if (pid < 0) {
		logger_log_error("master", "fork failed", errno);
		return -1;
	}
	if (pid == 0) {
		worker_loop(listen_fd);
		close(listen_fd);
		_exit(0);
	}
	return pid;
}
\end{lstlisting}

\begin{lstlisting}[caption={Функция создания слушающего сокета}, label={lst:listen}]
static int create_listen_socket(int port) {
	int fd = socket(AF_INET, SOCK_STREAM, 0);
	if (fd < 0) {
		logger_log_error("master", "socket failed", errno);
		return -1;
	}
	int reuse = 1;
	if (setsockopt(fd, SOL_SOCKET, SO_REUSEADDR, &reuse, sizeof(reuse)) < 0) {
		logger_log_error("master", "setsockopt(SO_REUSEADDR) failed", errno);
		close(fd);
		return -1;
	}
	struct sockaddr_in addr;
	memset(&addr, 0, sizeof(addr));
	addr.sin_family = AF_INET;
	addr.sin_addr.s_addr = htonl(INADDR_ANY);
	addr.sin_port = htons((uint16_t)port);
	if (bind(fd, (struct sockaddr *)&addr, sizeof(addr)) < 0) {
		logger_log_error("master", "bind failed", errno);
		close(fd);
		return -1;
	}
	if (listen(fd, LISTEN_BACKLOG) < 0) {
		logger_log_error("master", "listen failed", errno);
		close(fd);
		return -1;
	}
	return fd;
}
\end{lstlisting}

\begin{lstlisting}[caption={Функция работы worker'а}, label={lst:workerLoop}]
void worker_loop(int listen_fd) {
	printf("[worker %d] started\n", getpid());
	fflush(stdout);
	connection_t *conns = calloc(FD_SETSIZE, sizeof(connection_t));
	if (conns == NULL) {
		logger_log_error("worker", "calloc connections failed", errno);
		return;
	}
	for (int i = 0; i < FD_SETSIZE; ++i) {
		connection_init(&conns[i]);
	}
	if (set_nonblocking(listen_fd) < 0) {
		logger_log_error("worker", "failed to set listen socket nonblocking", errno);
	}
	int max_fd = listen_fd;
	while (!g_stop) {
		fd_set readfds;
		fd_set writefds;
		FD_ZERO(&readfds);
		FD_ZERO(&writefds);
		FD_SET(listen_fd, &readfds);
		for (int fd = 0; fd <= max_fd; ++fd) {
			if (conns[fd].active) {
				if (conns[fd].state == CONN_READING) {
					FD_SET(fd, &readfds);
				} else {
					FD_SET(fd, &writefds);
				}
			}
		}
		int ready = select(max_fd + 1, &readfds, &writefds, NULL, NULL);
		if (ready < 0) {
			if (errno == EINTR) {
				continue;
			}
			logger_log_error("worker", "select failed", errno);
			break;
		}
		if (FD_ISSET(listen_fd, &readfds)) {
			while (1) {
				struct sockaddr_storage client_addr;
				socklen_t client_len = sizeof(client_addr);
				int client_fd = accept(listen_fd, (struct sockaddr *)&client_addr, &client_len);
				if (client_fd < 0) {
					if (errno == EAGAIN || errno == EWOULDBLOCK) {
						break;
					}
					if (errno != EINTR) {
						logger_log_error("worker", "accept failed", errno);
					}
					break;
				}
				
				if (client_fd >= FD_SETSIZE) {
					close(client_fd);
					continue;
				}
				if (set_nonblocking(client_fd) < 0) {
					logger_log_error("worker", "failed to set client socket nonblocking", errno);
					close(client_fd);
					continue;
				}
				connection_init(&conns[client_fd]);
				conns[client_fd].active = true;
				conns[client_fd].state = CONN_READING;
				conns[client_fd].addr = client_addr;
				conns[client_fd].addr_len = client_len;
				if (client_fd > max_fd) {
					max_fd = client_fd;
				}
			}
		}
		for (int fd = 0; fd <= max_fd; ++fd) {
			if (!conns[fd].active) {
				continue;
			}
			if (conns[fd].state == CONN_READING && FD_ISSET(fd, &readfds)) {
				size_t left = sizeof(conns[fd].request_buf) - 1 - conns[fd].request_len;
				if (left == 0) {
					if (prepare_connection_response(&conns[fd]) < 0) {
						logger_log_error("worker", "prepare response failed for full request buffer",
						errno);
						close_connection(fd, &conns[fd]);
					}
					continue;
				}
				ssize_t n = recv(fd, conns[fd].request_buf + conns[fd].request_len, left, 0);
				if (n < 0) {
					if (errno == EAGAIN || errno == EWOULDBLOCK || errno == EINTR) {
						continue;
					}
					logger_log_error("worker", "recv failed", errno);
					close_connection(fd, &conns[fd]);
					continue;
				}
				if (n == 0) {
					close_connection(fd, &conns[fd]);
					continue;
				}
				conns[fd].request_len += (size_t)n;
				conns[fd].request_buf[conns[fd].request_len] = '\0';
				if (request_line_complete(&conns[fd])) {
					if (prepare_connection_response(&conns[fd]) < 0) {
						logger_log_error("worker", "http_prepare_response failed", errno);
						close_connection(fd, &conns[fd]);
					}
				}
			}
			if (conns[fd].active && conns[fd].state == CONN_WRITING && FD_ISSET(fd, &writefds)) {
				if (conns[fd].header_sent < conns[fd].plan.header_len) {
					const char *h = conns[fd].plan.header + conns[fd].header_sent;
					size_t len = conns[fd].plan.header_len - conns[fd].header_sent;
					ssize_t n = send(fd, h, len, MSG_NOSIGNAL);
					if (n < 0) {
						if (errno != EAGAIN && errno != EWOULDBLOCK && errno != EINTR) {
							logger_log_error("worker", "send header failed", errno);
							close_connection(fd, &conns[fd]);
						}
						continue;
					}
					conns[fd].header_sent += (size_t)n;
				}
				if (conns[fd].header_sent == conns[fd].plan.header_len && !conns[fd].body_done) {
					if (conns[fd].write_buf_sent == conns[fd].write_buf_len) {
						ssize_t nread = read(conns[fd].plan.file_fd, conns[fd].write_buf,
						sizeof(conns[fd].write_buf));
						if (nread < 0) {
							if (errno == EINTR) {
								continue;
							}
							logger_log_error("worker", "read file chunk failed", errno);
							close_connection(fd, &conns[fd]);
							continue;
						}
						if (nread == 0) {
							conns[fd].body_done = true;
						} else {
							conns[fd].write_buf_len = (size_t)nread;
							conns[fd].write_buf_sent = 0;
						}
					}
					if (!conns[fd].body_done && conns[fd].write_buf_sent < conns[fd].write_buf_len) {
						const char *p = conns[fd].write_buf + conns[fd].write_buf_sent;
						size_t len = conns[fd].write_buf_len - conns[fd].write_buf_sent;
						ssize_t n = send(fd, p, len, MSG_NOSIGNAL);
						if (n < 0) {
							if (errno != EAGAIN && errno != EWOULDBLOCK && errno != EINTR) {
								logger_log_error("worker", "send body chunk failed", errno);
								close_connection(fd, &conns[fd]);
							}
							continue;
						}
						conns[fd].write_buf_sent += (size_t)n;
						conns[fd].body_bytes_sent += n;
					}
				}
				if (conns[fd].active && conns[fd].header_sent == conns[fd].plan.header_len &&
				conns[fd].body_done) {
					finalize_and_close_connection(fd, &conns[fd]);
				}
			}
		}
		while (max_fd > listen_fd && !conns[max_fd].active) {
			--max_fd;
		}
	}
	for (int fd = 0; fd <= max_fd; ++fd) {
		if (conns[fd].active) {
			close_connection(fd, &conns[fd]);
		}
	}
	free(conns);
	printf("[worker %d] stopping\n", getpid());
	fflush(stdout);
}
\end{lstlisting}

\begin{lstlisting}[caption={Структура соединения}, label={lst:connection}]
typedef struct connection {
	bool active;
	connection_state_t state;
	struct sockaddr_storage addr;
	socklen_t addr_len;
	char request_buf[READ_BUFFER_SIZE];
	size_t request_len;
	http_response_plan_t plan;
	size_t header_sent;
	char write_buf[WRITE_BUFFER_SIZE];
	size_t write_buf_len;
	size_t write_buf_sent;
	bool body_done;
	bool timer_started;
	struct timespec started_at;
	long long body_bytes_sent;
} connection_t;
\end{lstlisting}
\clearpage
\begin{lstlisting}[caption={Структура плана HTTP-ответа}, label={lst:httpResponse}]
typedef struct http_response_plan {
	char method[16];
	char uri[2048];
	int status_code;
	char header[1024];
	size_t header_len;
	bool send_body;
	int file_fd;
	off_t file_size;
} http_response_plan_t;
\end{lstlisting}

\begin{lstlisting}[caption={Структура результата отправки HTTP-ответа}, label={lst:httpResponseResult}]
typedef struct http_result {
	char method[16];
	char uri[2048];
	int status_code;
	long long bytes_sent;
} http_result_t;
\end{lstlisting}

\begin{lstlisting}[caption={Функция подготовки плана ответа}, label={lst:responsePlan}]
int http_prepare_response(const char *raw, const char *docroot, http_response_plan_t *plan) {
	char method[32];
	char uri[2048];
	char version[32];
	char path[PATH_MAX];
	bool is_head = false;
	if (plan == NULL) {
		return -1;
	}
	plan_init(plan);
	if (parse_request_line(raw, method, sizeof(method), uri, sizeof(uri), version,
	sizeof(version)) < 0) {
		return build_error_header(plan, 400);
	}
	copy_trunc(plan->method, sizeof(plan->method), method);
	copy_trunc(plan->uri, sizeof(plan->uri), uri);
	if (strcmp(method, "GET") == 0) {
		is_head = false;
	} else if (strcmp(method, "HEAD") == 0) {
		is_head = true;
	} else {
		return build_error_header(plan, 405);
	}
	
	if (resolve_path(docroot, uri, path, sizeof(path)) < 0) {
		return build_error_header(plan, 403);
	}
	
	return prepare_file_response(path, is_head, plan);
}
\end{lstlisting}

\begin{lstlisting}[caption={Функция подготовки файла}, label={lst:responsePlanFile}]
static int prepare_file_response(const char *path, bool is_head, http_response_plan_t *plan) {
	struct stat st;
	if (stat(path, &st) < 0) {
		if (errno == EACCES) {
			return build_error_header(plan, 403);
		}
		return build_error_header(plan, 404);
	}
	if (!S_ISREG(st.st_mode)) {
		return build_error_header(plan, 403);
	}
	if (access(path, R_OK) < 0) {
		return build_error_header(plan, 403);
	}
	if (st.st_size < 0 || st.st_size > MAX_FILE_SIZE) {
		return build_error_header(plan, 403);
	}
	const char *content_type = guess_content_type(path);
	int n = snprintf(plan->header, sizeof(plan->header),
	"HTTP/1.1 200 OK\r\nConnection: close\r\nContent-Type: %s\r\n"
	"Content-Length: %lld\r\n\r\n",
	content_type, (long long)st.st_size);
	if (n <= 0 || (size_t)n >= sizeof(plan->header)) {
		return -1;
	}
	plan->header_len = (size_t)n;
	plan->status_code = 200;
	plan->file_size = st.st_size;
	if (is_head) {
		plan->send_body = false;
		plan->file_fd = -1;
		return 0;
	}
	int fd = open(path, O_RDONLY);
	if (fd < 0) {
		return build_error_header(plan, 403);
	}
	plan->file_fd = fd;
	plan->send_body = true;
	return 0;
}
\end{lstlisting}

\begin{lstlisting}[caption={Makefile для сборки и запуска сервера}, label={lst:make}]
CC := gcc
CFLAGS := -std=c11 -Wall -Wextra -Wpedantic -O2
LDFLAGS :=
TARGET := server
SRC := src/main.c src/worker.c src/http.c src/log.c
OBJ := $(SRC:.c=.o)
.PHONY: all clean run
all: $(TARGET)
$(TARGET): $(OBJ)
	$(CC) $(CFLAGS) -o $@ $^ $(LDFLAGS)
%.o: %.c
	$(CC) $(CFLAGS) -c -o $@ $<
run: $(TARGET)
	./$(TARGET)
clean:
	rm -f $(TARGET) $(OBJ)
\end{lstlisting}

\section{Демонстрация работы алгоритма}

По умолчанию сервер возвращает HTML-страницу c CSS-стилями. В листинге~\ref{lst:get} показан результат работы обработки запроса с методом \texttt{GET}, в листинге~\ref{lst:head}~---~запроса с методом \texttt{HEAD}, в листинге~\ref{lst:post}~---~запроса с не обрабатываемым методом \texttt{POST}.

\begin{lstlisting}[caption={Запрос GET}, label={lst:get}]
$ curl -i http://127.0.0.1:8080/
HTTP/1.1 200 OK
Connection: close
Content-Type: text/html; charset=utf-8
Content-Length: 424

<!doctype html>
<html lang="en">
<head>
	<meta charset="utf-8">
	<meta name="viewport" content="width=device-width, initial-scale=1">
	<title>Static Server</title>
	<link rel="stylesheet" href="/style.css">
</head>
<body>
	<main class="card">
		<h1>Static HTTP server</h1>
		<p>The server is running and serves files from docroot.</p>
		<p>Current MVP supports GET and HEAD methods.</p>
	</main>
</body>
</html>
\end{lstlisting}

\begin{lstlisting}[caption={Запрос HEAD}, label={lst:head}]
$ curl -I http://127.0.0.1:8080/
HTTP/1.1 200 OK
Connection: close
Content-Type: text/html; charset=utf-8
Content-Length: 424
\end{lstlisting}

\begin{lstlisting}[caption={Запрос HEAD}, label={lst:post}]
$ curl -i -X POST http://127.0.0.1:8080/
HTTP/1.1 405 Method Not Allowed
Connection: close
Allow: GET, HEAD
Content-Length: 0
\end{lstlisting}

\section{Тестирование}

В таблице~\ref{tab:tests} приведены тестовые сценарии и результаты тестирования разработанного HTTP-сервера.

\begin{sidewaystable}[]
	\centering
	\caption{Тестирование разработанного HTTP-сервера}
	\begin{tabular}{|p{5.5cm}|p{8cm}|p{5cm}|p{5cm}|}
		\hline
		Тест & Тестовый сценарий & Ожидаемый результат & Полученный результат\\\hline
		Запрос \texttt{GET} для получения файла по умолчанию & \$ curl -i http://127.0.0.1:8080/  & \texttt{200 OK}, файл в теле ответа & \texttt{200 OK}, файл в теле ответа \\\hline
		Запрос \texttt{GET} для получения файла & 
		\$ curl -i http://127.0.0.1:8080/manual\_test.txt & \texttt{200 OK}, файл в теле ответа & \texttt{200 OK}, файл в теле ответа \\\hline
		Запрос \texttt{HEAD}  & 
		\$ curl -I http://127.0.0.1:8080/manual\_test.txt & \texttt{200 OK} & \texttt{200 OK} \\\hline
		Запрос \texttt{GET} файла размером до 128 МБ  & 
		\$ curl -i http://127.0.0.1:8080/large\_120mb.txt & \texttt{200 OK}, файл в теле ответа & \texttt{200 OK}, файл в теле ответа \\\hline
		Запрос \texttt{POST} не поддерживаемого метода & \$ curl -i -X POST http://127.0.0.1:8080/ & \texttt{405 Method Not Allowed} & \texttt{405 Method Not Allowed} \\\hline
		Запрос \texttt{GET} файла размером свыше 128 МБ & 
		\$ curl -i http://127.0.0.1:8080/large\_129mb.txt & \texttt{403 Forbidden}, пустое тело ответа & \texttt{403 Forbidden}, пустое тело ответа \\\hline
		Запрос \texttt{GET} несуществующего файла  & 
		\$ curl -i http://127.0.0.1:8080/not\_found.txt & \texttt{404 Not Found}, пустое тело ответа & \texttt{404 Not Found}, пустое тело ответа \\\hline
		Запрос \texttt{GET} c path traversal & 
		\$ curl -i "http://127.0.0.1:8080/.{}./etc/passwd" & \texttt{403 Forbidden}, пустое тело ответа & \texttt{403 Forbidden}, пустое тело ответа \\\hline
		Запрос \texttt{GET} директории & \$ curl -i http://127.0.0.1:8080/dir & \texttt{403 Forbidden}, пустое тело ответа & \texttt{403 Forbidden}, пустое тело ответа \\\hline
		Запрос \texttt{GET} с невалидной стартовой строкой & \$ printf 'BROKEN\_REQUEST\verb|\r\n\r\n|' | nc 127.0.0.1 8080 & \texttt{400 Bad Request}, пустое тело ответа & \texttt{400 Bad Request}, пустое тело ответа \\\hline
	\end{tabular}
	\label{tab:tests}
\end{sidewaystable}

\section*{Вывод}

В данном разделе определены средства реализации, приведены реализации основных алгоритмов и функций, продемонстрирована работа HTTP-сервера и проведено его тестирование.
\clearpage